% ======================================================================
% Main file: The Σ-Model in Protein Domain Architectures
% NeurIPS 2027 submission
%
% FIXME: Swap \documentclass{neurips_2026} -> neurips_2027 when the
% official style file is released (expected Q1 2027). See SWAP-TO-2027.txt.
% ======================================================================
\documentclass[final]{neurips_2026}

% --- Fonts & encoding ---
\usepackage[utf8]{inputenc}
\usepackage[T1]{fontenc}
\usepackage{microtype}

% --- Core math & formatting ---
\usepackage{amsmath,amssymb,mathtools}
\usepackage{booktabs}
\usepackage{siunitx}
\usepackage{enumitem}

% --- Algorithms ---
\usepackage{algorithm}
\usepackage{algpseudocode}

% --- Figures ---
\usepackage{tikz}
\usepackage{pgfplots}
\pgfplotsset{compat=1.18}
\usetikzlibrary{patterns, positioning, arrows.meta, shapes.geometric}

% --- Custom commands ---
\newcommand{\R}{\mathbb{R}}
\DeclareMathOperator{\Var}{Var}
\DeclareMathOperator{\Corr}{Corr}
\DeclareMathOperator{\CosSim}{CosSim}

% ======================================================================
% Metadata
% ======================================================================
\title{The $\Sigma$-Model in Protein Domain Architectures:\\[2pt]
Schema-Coherence Suppression and Recombination}

\author{
  Basyirin Amsyar Basri \\
  Independent Researcher \\
  \texttt{basri@sigma-model.dev}
}

% ======================================================================
\begin{document}
% ======================================================================

\maketitle

\begin{abstract}
[PLACEHOLDER — 200-word abstract describing the application of the
$\Sigma$-Model framework to protein domain architecture recombination,
introducing the PfamCG-1.0 benchmark, and reporting experimental results
for brittle recombination, grammatical induction probes, and curriculum
design.]
\end{abstract}

\section{Introduction}
[PLACEHOLDER — Motivation: compositional generalisation in protein
domain architectures. The $\Sigma$-Model \cite{basri2026sigma} as a
dynamical-systems framework for schema-coherence suppression. Three
contributions: brittle recombination (N=500), grammatical probes
(3 probe types), and PfamCG-1.0 benchmark suite.]

\section{Background: The $\Sigma$-Model Framework}
[PLACEHOLDER — Summary of the coupled ODE system
(Eqs.~1--6 of Basri, 2026), key concepts ($\sigma_A$, $\delta_A$,
$\alpha_A$), phases of learning, and the training-aware curriculum
hypothesis.]

\section{Pfam Domain Architectures as a Test Bed}
[PLACEHOLDER — Pfam database (version 39.0), domain architecture
representation, recombination distance via HMM bit scores, vocabulary
construction, and dataset statistics.]

\section{Experiment 1: Brittle Domain Recombination}
[PLACEHOLDER — N=500 regime: familiar combinations, new families, and
cross-fold transfer. Hypothesis: $\Sigma$-Model predicts brittle
recombination for unfamiliar domain pairs. Results: accuracy vs.\
familiarity, effect size, statistical tests.]
% CLAIM:C-039 — σ-trap replicates in protein domain data (H1)
% CLAIM:C-040 — ΔCG ≥ 30pp for low-σ models on Pfam OOD splits

\section{Experiment 2: Grammatical Induction Probes}
[PLACEHOLDER — Three probe types: positional bias (linear),
co-occurrence (bilinear), hierarchy (MLP). Results per probe:
accuracy, $R^2$, and comparison to random baseline.]
% CLAIM:C-041 — PLMs encode surface co-occurrence, not domain grammar (H2)
% CLAIM:C-042 — Structural OOD accuracy < lexical OOD accuracy by >10pp

\section{Experiment 3: PfamCG-1.0 Benchmark}
[PLACEHOLDER — Benchmark results: recombination accuracy, domain
retention rate, novel combination precision, fold consistency.
Comparison with baselines: random, frequency-based, co-occurrence.]
% CLAIM:C-043 — Multiplicative model fits better than additive (H3)
% CLAIM:C-044 — PfamCG-1.0 benchmark valid against known biology

\section{Discussion}
[PLACEHOLDER — Implications for protein engineering, limitations
(vocabulary size, architecture depth), connections to other
compositional generalisation benchmarks.]

\section{Related Work}
[PLACEHOLDER — Protein language models (ESM, ProtBERT), domain
architecture prediction, curriculum learning for biological sequences,
compositional generalisation in NLP.]

\section{Conclusion}
[PLACEHOLDER — Summary of findings and future directions.]

% ======================================================================
\bibliographystyle{plain}
\bibliography{bibliography}
% ======================================================================

\appendix
\section{Supplemental Materials}
[PLACEHOLDER — Additional experimental details, hyperparameter grids,
and full result tables.]

\end{document}
